\section{Conclusion}
\label{sec:conclusion}

\paragraph{Summary.}
We have presented three studies of pretrained ImageNet networks that,
together with two honest negatives, witness a single thesis: the
knowledge matrix is a \emph{canonical arrangement} of a trained network's
per-sample behaviour that penultimate-layer activations cannot supply.
Study~1 (Section~\ref{sec:study1}) is a correctness-and-measurement
study. Its permutation arm verifies the permutation invariance of
Proposition~\ref{prop:perm} on ResNet-152 --- where invariance is automatic
by construction, so this is a numerical correctness check rather than
independent evidence (the load-bearing, non-automatic invariance of
Theorem~\ref{thm:maxinv} is the cross-architecture case): measured signal-relative, random neuron
permutations leave the knowledge matrix with a permutation drift of only
$0.1$--$0.2\%$ of its own adversarial-pair signal (the residual being
float32 accumulation, $\sim10^{-10}$ at float64), while the penultimate
features move by an amount of the same order as their own adversarial
signal. Its teleportation arm measures, rather than asserts, the
knowledge-matrix drift across all three architectures (ResNet-152,
DenseNet-121, GoogLeNet): teleportation visibly perturbs the penultimate
features on every architecture (the similarity panel of
Section~\ref{sec:study1}), while the knowledge matrix does not: the transform is
an exact isomorphism, batch normalisation included, and the matrix's drift under
it is measured at $0.8$--$2.5\times10^{-15}$ relative --- the double-precision
floor --- on all three architectures (Appendix~\ref{app:teleport-exact}).
Study~2 (Section~\ref{sec:study2}) is descriptive geometry: it applies
the visible/invisible decomposition of Theorem~\ref{thm:pyth} to
adversarial pairs and reads the resulting coherence $A=(d_\Psi/d_M)^2$
against the single-pixel line $A=1$. On the deep ImageNet trio every
attack family moves the germ with median $A\le0.23$, well below that
line, and the attack-family ordering of $A$ is stable across six
rater architectures (Kendall $W=0.921$), the population-matched full-scale
appendix panel reproducing the same family-level grouping on all three of
its architectures (Table~\ref{tab:s3_ordering_panel}). Study~3
(Section~\ref{sec:study4}) is an availability result: the fixed-shape
knowledge matrix lets one compare architectures directly, with no
alignment step. The two honest negatives (Section~\ref{sec:study3}) are
the detector bake-off, in which penultimate features win $5$ of $6$
configurations and the original detection claim is withdrawn, and an
LP-counterfactual that fails structurally on pretrained ImageNet networks
(\texttt{region\_ok} = 0/54, Limitation L4).

\paragraph{The umbrella thesis, restated.}
Hidden activations are gauge-covariant and germ-incomplete
(Theorem~\ref{thm:complete}): they change under
network isomorphisms that preserve the function exactly, and they
admit no exact accounting that ties their displacement to the change
in the network's output. The knowledge matrix avoids both: it is an
invariant of the function (Theorem~\ref{thm:maxinv}), and every
displacement splits exactly into a logit-visible and a logit-invisible
part whose balance is the unit-free coherence $A$
(Theorem~\ref{thm:pyth}). For the piecewise-linear networks studied
here, these function-level properties --- invariance and the exact
completeness identity --- are shared with per-class gradient$\times$input
plus an aggregate bias attribution (Theorem~\ref{thm:weff}); what the
knowledge matrix contributes is the canonical \emph{arrangement} of that
data, which is what makes the alignment-free cross-architecture
comparison possible and what the $C$-VJP computation of
Appendix~\ref{app:engineering} exploits. The three studies above, with
their two negatives, are the first witnesses of this thesis.

\paragraph{HAaNE~I is foundational; HAaNE~II asks a different question.}
This paper is HAaNE~I of an ongoing series, and its standalone
contribution is deliberately foundational: a canonical per-sample object,
a characterisation of \emph{exactly} which transformations leave it
invariant (the full function-stabiliser, Theorem~\ref{thm:maxinv}, of
which neuron permutation is one special case), and the resulting
alignment-free cross-architecture comparison that penultimate features
cannot support. We make no application claim in this paper; the
applications that motivated the original framing belong to a separate,
unnumbered companion. Whether independently trained networks share the
knowledge matrix, and how that depends on training, is the subject of
HAaNE~II.

\paragraph{Future directions.}
Subsequent installments take up different questions about the same
object. The next paper, HAaNE~II, asks whether independently trained
networks learn the same knowledge matrix, and how the answer depends on
width, learning rate and parameterisation. Several concrete follow-ups
suggest themselves from the Limitations of this paper. Constructing two
genuinely different architectures that realise the same function on a
neighbourhood --- by dead-unit insertion, neuron splitting, or a width-changing
re-encoding --- and checking that their knowledge matrices agree would supply the
transform-based evidence for cross-architecture invariance that no
same-architecture symmetry can give (closes L1).
A multi-region or continuation reformulation of the
$\M(x)$-counterfactual would convert the structural negative of
Section~\ref{sec:study3} into a positive one (closes L4). Concrete
experiments demonstrating the practical consequences of the
canonical-arrangement claim --- federated-learning aggregation in matrix
space, model comparison without CKA-style alignment, $k$-nearest-neighbour
classification on knowledge matrices --- are the natural application track
of the separate, unnumbered companion (closes L5). Cross-pollination with
transformer-interpretability methods --- training sparse autoencoders on
knowledge-matrix rows
\citep{bricken2023monosemanticity,templeton2024scaling}, projecting
through partial row sums in analogy with the logit lens
\citep{nostalgebraist2020logitlens} --- would extend the framework
beyond the CNN setting of this paper (closes L6).

\paragraph{Closing.}
Hidden activations are not enough. A trained network's behaviour
on a sample is a property of its \emph{function}, and --- unlike the
penultimate activations --- the knowledge matrix derived from its
quiver representation is an invariant of that function
(Theorem~\ref{thm:maxinv}) with an exact accounting of how it moves.
The studies in this paper are the first three
witnesses in a longer multi-witness argument.
